\documentclass[11pt]{article}
\usepackage[margin=0.85in]{geometry}
\usepackage{fontspec}
\usepackage{titlesec}
\usepackage{enumitem}
\usepackage{xcolor}
\usepackage[skip=4pt]{parskip}
\setlist{itemsep=1pt, topsep=1pt}

\setmainfont{Calibri}[
  UprightFont = *,
  BoldFont = * Bold,
  ItalicFont = * Italic
]

\definecolor{accent}{HTML}{1F4E79}

\titleformat{\section}{\large\bfseries\color{accent}}{\thesection.}{0.5em}{}
\titlespacing{\section}{0pt}{0.5em}{0.3em}

\setlength{\parindent}{0pt}

\begin{document}

\begin{center}
    {\LARGE\bfseries\color{accent} Sampling Plan: ADHD Prevalence in Aguascalientes}\\[0.4em]
    {\large Homework 2: Sampling}
\end{center}

\vspace{0.5em}

\begin{center}
\begin{tabular}{ll}
\textbf{Name:} & Alexis Alberto Z\'u\~niga Alonso \\
\textbf{Institution:} & EdgeHub | School of Innovation \\
\textbf{Class:} & Statistical Methods \\
\textbf{Submitted to:} & M.Sc. Axel Torres \\
\textbf{Date:} & September 25, 2026 \\
\end{tabular}
\end{center}

\vspace{0.4em}
\hrule
\vspace{0.2em}

\section{Problem Statement}
This study looks at how common ADHD is among children and teenagers in the state of Aguascalientes, Mexico. The goal is to get a reliable estimate of ADHD screening rates in school-age students, which can help schools and health authorities plan better support and resources.

\section{Population}
All children and teenagers between 6 and 17 years of age enrolled in elementary and high schools in the state of Aguascalientes.

\section{Variable to Measure}
Whether the student screens positive for ADHD (yes or no). This is a categorical, binary variable.

\section{Unit of Observation}
Each individual student between 6 and 17 years old, enrolled in a school in Aguascalientes, on whom an ADHD screening result (positive or negative) is recorded.

\section{Sampling Frame}
The \emph{Cat\'alogo de Centros de Trabajo de la SEP}, which lists every public and private elementary and high school currently active in Aguascalientes. Once a school is selected, its own enrollment lists per group are used as the frame for the second stage.

\section{Sampling Unit}
Two sampling units are used, one for each stage of the design:
\begin{itemize}[itemsep=2pt]
    \item \textbf{First stage:} the school. Each school registered in the SEP catalog is a sampling unit.
    \item \textbf{Second stage:} the group (classroom). Within each school selected in the first stage, groups are the sampling units from which students are drawn.
\end{itemize}

\section{Sampling Method}
Probabilistic sampling. Every school and every group has a known, nonzero chance of being selected, which allows the results to be generalized to the full population and lets us calculate a margin of error.

\section{Type of Sampling}
Two-stage stratified sampling.
\begin{itemize}[itemsep=2pt]
    \item \textbf{Stage 1 (schools):} Schools are first grouped into strata, for example by municipality and by school type (public or private). A random sample of schools is then drawn from each stratum.
    \item \textbf{Stage 2 (groups):} Within each selected school, groups are randomly sampled, and all students in the chosen groups are screened for ADHD.
\end{itemize}
This design keeps costs and travel down while still giving fair representation to the different types of schools in the state.

\end{document}
